% GENERATED FILE. Edit canonical lessons, then run npm run generate:books.
% canonical-source-hash: fnv1a64:ba28f0aa9519c6fc
% canonical-lessons: GU-C329-rangoli, GU-C329-sangarvu, GU-C329-jhagmagvu, GU-C329-mala, GU-C329-sobhayatra

\chapter{Lights and Colours}
\label{ch:gu-lights-and-colours}
\hlchaptermodality{329}

\begin{chapteropening}
\textbf{By the end of this chapter:} \emph{I can say how a town is made ready for a festival.}

Complete the last lesson of chapter 329: I can say how a town is made ready for a festival.
\end{chapteropening}

\section[\texorpdfstring{\gu{રંગોળી} (rangoḷī)}{રંગોળી (rangoḷī)}]{\gu{રંગોળી} (rangoḷī) — a rangoli (a coloured floor pattern)}

\label{lesson:GU-C329-rangoli}



\begin{quote}
Before the new one: say the Gujarati for to light, then the Gujarati for incense.
\end{quote}

\subsection*{You'll want to know: \gu{રંગોળી}}
\textbf{\gu{રંગોળી}} — \emph{rangoḷī} — \textquotedblleft{}a rangoli (a coloured floor pattern)\textquotedblright{}.

\begin{quote}
\gu{દિવાળીમાં} \gu{મીરાએ} \gu{આંગણામાં} \gu{રંગોળી} \gu{બનાવી}.
\end{quote}

\begin{quote}
\emph{divāḷīmā̃ Mīrāe ā̃gṇāmā̃ rangoḷī banāvī.}
\end{quote}

\begin{quote}
At Diwali Mira made a rangoli in the courtyard.
\end{quote}

A feminine noun.

\begin{cousinweb}
\textbf{\gu{રંગોળી}} (\emph{rangoḷī}) has a history that is not settled: Monier-Williams records a Sanskrit \emph{raṅgavallī}, built from \emph{raṅga}, colour, and \emph{vallī}, a creeper, but only as the name of a plant used at sacrifices. It is a pattern of coloured powder laid on the floor at a doorway for a festival; it is feminine: \textbf{\gu{રંગોળી} \gu{બનાવી}}.
\end{cousinweb}

\begin{grammarlens}[title={What you've built}]
One more word for telling what happened.
\end{grammarlens}

\subsection*{Your turn}
\begin{itemize}
\raggedright
  \item \emph{Say it:} \emph{rangoḷī}
  \item \emph{Say it:} \emph{rangoḷī}, once more
  \item \emph{Say it:} \emph{divāḷīmā̃ Mīrāe ā̃gṇāmā̃ rangoḷī banāvī.} — At Diwali Mira made a rangoli in the courtyard.
  \item \emph{Recall:} say the Gujarati for to light, then the Gujarati for incense, then say \emph{rangoḷī} again
\end{itemize}

\begin{tcolorbox}[breakable,colback=teal!4,colframe=teal!35!black,title={Before you move on}]
What does \gu{રંગોળી} mean? (\textquotedblleft{}A rangoli (a coloured floor pattern)\textquotedblright{}.) Say it once more.
\end{tcolorbox}
\section[\texorpdfstring{\gu{શણગારવું} (śaṇgārvũ)}{શણગારવું (śaṇgārvũ)}]{\gu{શણગારવું} (śaṇgārvũ) — to adorn}

\label{lesson:GU-C329-sangarvu}



\begin{quote}
Before the new one: say the Gujarati for incense, then the Gujarati for a rangoli.
\end{quote}

\subsection*{You'll want to know: \gu{શણગારવું}}
\textbf{\gu{શણગારવું}} — \emph{śaṇgārvũ} — \textquotedblleft{}to adorn\textquotedblright{}.

\begin{quote}
\gu{લોકોએ} \gu{આખું} \gu{ગામ} \gu{શણગાર્યું}.
\end{quote}

\begin{quote}
\emph{lokoe ākhũ gām śaṇgāryũ.}
\end{quote}

\begin{quote}
People decorated the whole village.
\end{quote}

Past, with \gu{તેણે} (he or she …): \gu{તેણે} \gu{ઘર} \gu{શણગાર્યું}, \emph{teṇe ghar śaṇgāryũ}.

\begin{cousinweb}
comes from \emph{*śṛṅgārayati}, \textquotedblleft{}adorns\textquotedblright{}, a form CDIAL reconstructs on \emph{śṛṅgāra}, finery (CDIAL 12593). \textbf{\gu{સજાવવું}} is to decorate; \textbf{\gu{શણગારવું}} is to dress up richly, a bride or a whole town.
\end{cousinweb}

\begin{grammarlens}[title={What you've built}]
One more word for telling what happened.
\end{grammarlens}

\subsection*{Your turn}
\begin{itemize}
\raggedright
  \item \emph{Say it:} \emph{śaṇgārvũ}
  \item \emph{Say it:} \emph{śaṇgārvũ}, once more
  \item \emph{Say it:} \emph{lokoe ākhũ gām śaṇgāryũ.} — People decorated the whole village.
  \item \emph{Recall:} say the Gujarati for incense, then the Gujarati for a rangoli, then say \emph{śaṇgārvũ} again
\end{itemize}

\begin{tcolorbox}[breakable,colback=teal!4,colframe=teal!35!black,title={Before you move on}]
What does \gu{શણગારવું} mean? (\textquotedblleft{}To adorn\textquotedblright{}.) Say it once more.
\end{tcolorbox}
\section[\texorpdfstring{\gu{ઝગમગવું} (jhagmagvũ)}{ઝગમગવું (jhagmagvũ)}]{\gu{ઝગમગવું} (jhagmagvũ) — to sparkle}

\label{lesson:GU-C329-jhagmagvu}



\begin{quote}
Before the new one: say the Gujarati for a rangoli, then the Gujarati for to adorn.
\end{quote}

\subsection*{You'll want to know: \gu{ઝગમગવું}}
\textbf{\gu{ઝગમગવું}} — \emph{jhagmagvũ} — \textquotedblleft{}to sparkle\textquotedblright{}.

\begin{quote}
\gu{દિવાળીની} \gu{રાત્રે} \gu{આખું} \gu{શહેર} \gu{દીવાથી} \gu{ઝગમગતું} \gu{હતું}.
\end{quote}

\begin{quote}
\emph{divāḷīnī rātre ākhũ śaher dīvāthī jhagmagtũ hatũ.}
\end{quote}

\begin{quote}
On Diwali night the whole city sparkled with lamps.
\end{quote}

Past (it …): \gu{તે} \gu{શહેર} \gu{ઝગમગ્યું}, \emph{te śaher jhagmagyũ}.

\begin{cousinweb}
\textbf{\gu{ઝગમગવું}} (\emph{jhagmagvũ}) is an imitative word: CDIAL sets the Hindi cousin \emph{jagmag}, \textquotedblleft{}glitter\textquotedblright{}, under \emph{*jag}, \textquotedblleft{}glitter\textquotedblright{}, a form it reconstructs (CDIAL 5076), but does not list the Gujarati verb. The doubled sound \emph{jhag-mag} imitates light that flickers and shines.
\end{cousinweb}

\begin{grammarlens}[title={What you've built}]
One more word for telling what happened.
\end{grammarlens}

\subsection*{Your turn}
\begin{itemize}
\raggedright
  \item \emph{Say it:} \emph{jhagmagvũ}
  \item \emph{Say it:} \emph{jhagmagvũ}, once more
  \item \emph{Say it:} \emph{divāḷīnī rātre ākhũ śaher dīvāthī jhagmagtũ hatũ.} — On Diwali night the whole city sparkled with lamps.
  \item \emph{Recall:} say the Gujarati for a rangoli, then the Gujarati for to adorn, then say \emph{jhagmagvũ} again
\end{itemize}

\begin{tcolorbox}[breakable,colback=teal!4,colframe=teal!35!black,title={Before you move on}]
What does \gu{ઝગમગવું} mean? (\textquotedblleft{}To sparkle\textquotedblright{}.) Say it once more.
\end{tcolorbox}
\section[\texorpdfstring{\gu{માળા} (māḷā)}{માળા (māḷā)}]{\gu{માળા} (māḷā) — a garland}

\label{lesson:GU-C329-mala}



\begin{quote}
Before the new one: say the Gujarati for to adorn, then the Gujarati for to sparkle.
\end{quote}

\subsection*{You'll want to know: \gu{માળા}}
\textbf{\gu{માળા}} — \emph{māḷā} — \textquotedblleft{}a garland\textquotedblright{}.

\begin{quote}
\gu{મીરાએ} \gu{ફૂલની} \gu{માળા} \gu{બનાવી}.
\end{quote}

\begin{quote}
\emph{Mīrāe phūlnī māḷā banāvī.}
\end{quote}

\begin{quote}
Mira made a garland of flowers.
\end{quote}

A feminine noun.

\begin{cousinweb}
\textbf{\gu{માળા}} (\emph{māḷā}) is a tatsama, Sanskrit \emph{mālā}, \textquotedblleft{}a wreath, a garland\textquotedblright{} (CDIAL 10092). It is feminine: \textbf{\gu{ફૂલની} \gu{માળા}}.
\end{cousinweb}

\begin{grammarlens}[title={What you've built}]
One more word for telling what happened.
\end{grammarlens}

\subsection*{Your turn}
\begin{itemize}
\raggedright
  \item \emph{Say it:} \emph{māḷā}
  \item \emph{Say it:} \emph{māḷā}, once more
  \item \emph{Say it:} \emph{Mīrāe phūlnī māḷā banāvī.} — Mira made a garland of flowers.
  \item \emph{Recall:} say the Gujarati for to adorn, then the Gujarati for to sparkle, then say \emph{māḷā} again
\end{itemize}

\begin{tcolorbox}[breakable,colback=teal!4,colframe=teal!35!black,title={Before you move on}]
What does \gu{માળા} mean? (\textquotedblleft{}A garland\textquotedblright{}.) Say it once more.
\end{tcolorbox}
\section[\texorpdfstring{\gu{શોભાયાત્રા} (śobhāyātrā)}{શોભાયાત્રા (śobhāyātrā)}]{\gu{શોભાયાત્રા} (śobhāyātrā) — a procession}

\label{lesson:GU-C329-sobhayatra}



\begin{quote}
Before the new one: say the Gujarati for to sparkle, then the Gujarati for a garland.
\end{quote}

\subsection*{You'll want to know: \gu{શોભાયાત્રા}}
\textbf{\gu{શોભાયાત્રા}} — \emph{śobhāyātrā} — \textquotedblleft{}a procession\textquotedblright{}.

\begin{quote}
\gu{સાંજે} \gu{ગામમાં} \gu{મોટી} \gu{શોભાયાત્રા} \gu{નીકળી}.
\end{quote}

\begin{quote}
\emph{sā̃je gāmmā̃ moṭī śobhāyātrā nīkḷī.}
\end{quote}

\begin{quote}
In the evening a big procession went through the village.
\end{quote}

A feminine noun.

\begin{cousinweb}
\textbf{\gu{શોભાયાત્રા}} (\emph{śobhāyātrā}) is built from Sanskrit \emph{śobhā}, \textquotedblleft{}splendour\textquotedblright{}, and \emph{yātrā}, \textquotedblleft{}a going, a journey\textquotedblright{} (Monier-Williams), the word behind \textbf{\gu{યાત્રા}}: a splendid going through the streets. A procession \textbf{\gu{નીકળે} \gu{છે}}, sets out.
\end{cousinweb}

\begin{grammarlens}[title={What you've built}]
One more word for telling what happened.
\end{grammarlens}

\subsection*{Your turn}
\begin{itemize}
\raggedright
  \item \emph{Say it:} \emph{śobhāyātrā}
  \item \emph{Say it:} \emph{śobhāyātrā}, once more
  \item \emph{Say it:} \emph{sā̃je gāmmā̃ moṭī śobhāyātrā nīkḷī.} — In the evening a big procession went through the village.
  \item \emph{Recall:} say the Gujarati for to sparkle, then the Gujarati for a garland, then say \emph{śobhāyātrā} again
\end{itemize}

\begin{tcolorbox}[breakable,colback=teal!4,colframe=teal!35!black,title={Before you move on}]
What does \gu{શોભાયાત્રા} mean? (\textquotedblleft{}A procession\textquotedblright{}.) Say it once more.
\end{tcolorbox}
